\begin{textsample}{POS Dim 1 – 2015\_c4 – Score -3.00 – 2015\_c4/enade\_2015\_tec\_log\_3a04996ad2.txt}  \label{ex:f1_pos_332}
CAPÍTULO 4 PERCEPÇÃO DA PROVA As análises feitas neste capítulo tratam das percepções dos concluintes da Área de Tecnologia em Logística sobre a prova aplicada no Enade/2015.
Essas percepções foram mensuradas por meio de nove questões que avaliaram desde o grau de dificuldade da prova até o tempo gasto para concluí -la.
As percepções sobre a prova foram relacionadas com o desempenho dos estudantes e com a Grande Região de funcionamento do curso.
O questionário de percepção da prova encontra -se ao final do Anexo VII, que traz a reprodução da prova.
O desempenho dos estudantes foi classificado em quatro quartos.
Para tanto, esse desempenho foi ordenado de forma ascendente.
O percentil 25, P25, também conhecido como primeiro quartil, é a nota de desempenho que deixa um quarto ( 25 \% ) dos valores observados abaixo e três quartos acima.
A Figura 1 apresenta uma ilustração deste conceito.
O quarto inferior de desempenho é composto pelas notas abaixo do primeiro quartil.
Já o percentil 75, P75, também conhecido como terceiro quartil, é o valor para o qual há três quartos ( 75 \% ) dos dados abaixo e um quarto acima dele.
O quarto superior de desempenho é composto pelas notas iguais ou acima do terceiro quartil.
O percentil 50, P50, também conhecido como mediana, é o valor que divide as notas em dois conjuntos de igual tamanho.
O segundo quarto inclui valores entre o primeiro quartil ( P25 ) e a mediana.
O terceiro quarto contém os valores entre a mediana ( P50 ) e o terceiro quartil ( P75 ).
Vale ressaltar que percentis, quartis e medianas são pontos que não obrigatoriamente pertencem ao conjunto original de dados, ao passo que os quartos são subconjuntos dos dados originais.
A seguir, serão apresentados gráficos com resultados selecionados, relativos às nove questões avaliadas por grupos de estudantes.
Os gráficos apresentam nas barras o percentual de alunos que assinalaram uma das opções ou a soma das porcentagens daqueles que assinalaram duas ( ou três ) delas.
Por exemplo, para as questões 1 e 2, os gráficos apresentam a porcentagem total de participantes que assinalaram as opções ( D ) Difícil e ( E ) Muito difícil.
Em cada barra foram assinalados também os extremos do intervalo de confiança de 95 \% como linhas verticais unidas por uma linha horizontal na forma da letra H maiúscula.
As Tabelas no Anexo II apresentam os valores absolutos e a distribuição percentual das alternativas válidas das nove questões, segundo o mesmo recorte de desempenho dos alunos e Grande Região de funcionamento do curso. 4.1 GRAU DE DIFICULDADE DA PROVA 4.1.1 Componente de Formação Geral Ao avaliarem “ Qual o grau de dificuldade desta prova na parte de Formação Geral? ” ( Questão 1 ), 33,1 \% do grupo de inscritos e presentes optaram pelas alternativas Difícil ou Muito difícil.
Entretanto, para mais da metade dos estudantes ( 58,6 \% ), o Componente de Formação Geral da prova foi considerado com grau de dificuldade Médio ( Gráfico 4.1, Gráfico 4.2 e, no Anexo II, a Tabela II.1 ).
O percentual de estudantes que consideraram a prova como Difícil ou Muito difícil foi maior na região Sul, onde a proporção foi de 35,4 \%, enquanto a de menor incidência foi a Nordeste, com 24,5 \%, observando -se diferenças estatisticamente significativas dessa região para as demais, conforme o Gráfico 4.1.
Nas Grandes Regiões, a proporção de presentes à prova que consideraram o Componente de Formação Geral como sendo de grau de dificuldade Médio esteve entre 56,6 \% na região Sul e 63,1 \% na região Nordeste.
O percentual de alunos que consideraram a prova Difícil ou Muito difícil foi decrescente em função dos quartos de desempenho : 36,6 \% no primeiro quarto e 29,1 \% no quarto, grupo de melhor desempenho na prova.
Nos quartos de desempenho intermediários, a proporção de alunos que consideraram a prova Difícil ou Muito difícil foi de 35,2 \% no segundo quarto e 31,8 \% no terceiro quarto.
Há diferenças estatisticamente significativas do primeiro quarto para o terceiro e do primeiro e segundo para o quarto superior.
Para todos os quartos de desempenho, a alternativa modal para esta pergunta foi Médio, com 56,6 \% e 60,0 \% dos respondentes nos quartos extremos, primeiro e quarto, respectivamente. 4.1.2 Componente de Conhecimento Específico Ao responderem à Questão 2 – “ Qual o grau de dificuldade desta prova na parte de Componente Específico? ” – 38,9 \% do grupo de estudantes classificaram -na como Difícil ou Muito difícil.
Além disso, o Componente de Conhecimento Específico da prova foi considerado com grau de dificuldade Médio por 56,0 \% dos alunos ( Gráfico 4.3, Gráfico 4.4, e, no Anexo II, a Tabela II.2 ).
A análise das respostas dos estudantes quanto ao grau de dificuldade do Componente de Conhecimento Específico da prova, agregado por Grande Região, mostra que a diferença entre a maior e a menor proporção de alunos que a avaliaram como Difícil ou Muito difícil é estatisticamente significativa : a maior na região Norte ( 42,3 \% ) e a menor proporção na região Nordeste ( 27,3 \% ).
Também ocorre diferença estatisticamente significativa entre a região Nordeste e as regiões Sudeste ( 39,9 \% ) e Sul ( 41,2 \% ).
O percentual de alunos que classificaram o grau de dificuldade como Médio, no Componente de Conhecimento Específico, variou de 52,0 \% a 64,9 \%, para as regiões Norte e Nordeste, respectivamente.
Considerando -se a avaliação da dificuldade das questões do Componente de Conhecimento Específico da prova, de acordo com o desempenho dos estudantes, verifica -se que a proporção dos que classificaram a parte específica como Difícil ou Muito difícil variou de 38,1 \% ( último quarto ) a 40,1 \% ( primeiro quarto ), não sendo observadas diferenças estatisticamente significativas do resultado entre os quatro quartos de desempenho.
A alternativa modal para a Questão 2 foi o grau médio, com 54,4 \% do quarto inferior e 57,1 \% do superior optando por essa resposta. 4.2 EXTENSÃO DA PROVA EM RELAÇÃO AO TEMPO TOTAL Indagados quanto à extensão da prova, em relação ao tempo total oferecido para a sua resolução ( Questão 3 ), os estudantes apontaram, com maior incidência, a alternativa que considerava a extensão adequada, para todas as agregações consideradas ( Gráfico 4.5, Gráfico 4.6, e, no Anexo II, a Tabela II.3 ).
O percentual de alunos que responderam ser a extensão da prova adequada foi de 58,3 \%.
Já 34,7 \% dos inscritos presentes consideraram que a prova foi longa ou muito longa e 7,0 \% a avaliaram como curta ou muito curta.
Entre as Grandes Regiões a proporção daqueles que avaliaram a prova como longa ou muito longa em relação ao tempo total destinado à sua resolução variou pouco : de 30,5 \% na região Centro-Oeste até 36,0 \% na região Sudeste.
As diferenças entre a região Sudeste e as regiões Norte ( 30,7 \% ) e Nordeste ( 30,8 \% ) são estatisticamente significativas.
A escolha modal entre as grandes regiões foi que a extensão da prova foi Adequada, variando de 56,2 \% ( Nordeste ) a 59,8 \% ( Sul ).
Considerando -se o desempenho dos alunos, nota -se ainda que 57,9 \% consideraram a extensão da prova adequada no quarto de desempenho inferior e 59,7 \% no de melhor desempenho ( quarto superior ).
Nos quartos intermediários, essa proporção foi 58,0 \% no segundo quarto e 57,6 \% no terceiro.
No Gráfico 4.6, observam -se as proporções de estudantes que consideraram a prova longa ou muito longa em função dos quartos de desempenho, com variação entre 33,7 \% ( quarto superior ) e 35,4 ( terceiro quarto ), não sendo observadas diferenças estatisticamente significativas entre as proporções dos quatro quartos de desempenho. 4.3 COMPREENSÃO DOS ENUNCIADOS DAS QUESTÕES 4.3.1 Componente de Formação Geral Com relação aos enunciados das questões do Componente de Formação Geral ( Questão 4 ), as opiniões foram positivas, já que 73,3 \% dos alunos avaliados consideraram todos ou a maioria dos enunciados das questões “ claros e objetivos ” ( Gráfico 4.7, Gráfico 4.8, e, no Anexo II, a Tabela II.4 ).
Na análise regional, a percentagem de estudantes que avaliaram que todos ou a maioria dos enunciados das questões do Componente de Formação Geral estavam “ claros e objetivos ” variou de 72,1 \% na região Sudeste a 79,8 \% na região Centro-Oeste, sendo essa a única diferença estatisticamente significativa.
A análise das percepções dos estudantes sobre a clareza e objetividade dos enunciados permite afirmar que todos, ou a maioria dos enunciados de questões relativas ao Componente de Formação Geral, foram considerados “ claros e objetivos ” para a maior parte dos respondentes ( maior do que 72 \% em todas as regiões e maior do que 65 \% para todos os quartos de desempenho ).
Segundo o desempenho, observa -se que a proporção dos que emitiram essa opinião cresce conforme o desempenho aumenta, com diferenças estatisticamente 94 significativas entre todos os quartos.
No quarto superior, a clareza e objetividade de todos ou da maioria dos enunciados das questões foi percebida por 81,2 \% dos alunos.
Para os demais, as proporções foram 75,5 \% ( terceiro quarto ), 70,9 \% ( segundo quarto ) e 65,1 \% ( quarto inferior ). 4.3.2 Componente de Conhecimento Específico Com relação aos enunciados das questões do Componente de Conhecimento Específico da prova, para 74,2 \% dos estudantes avaliados da Área de Tecnologia em Logística, a clareza e a objetividade ( Questão 5 ) estavam presentes em todas ou na maioria das questões ( Gráfico 4.9, Gráfico 4.10, e no Anexo II, a Tabela II.5 ).
A maioria dos estudantes de todas as Grandes Regiões brasileiras considerou “ claros e objetivos ” todos ou a maioria dos enunciados das questões do Componente de Conhecimento Específico da prova, percentual sempre maior que 72 \%.
A maior proporção foi na região Centro-Oeste ( 79,2 \% ), e a menor na região Norte ( 72,3 \% ).
Observam -se diferenças estatisticamente significativas entre a região Sudeste ( 73,0 \% ) e as regiões Sul ( 76,3 \% ) e Centro-Oeste ( 79,2 \% ).
A proporção de estudantes que consideraram os enunciados das questões “ claros e objetivos ” apresenta uma tendência crescente em relação ao aumento de desempenho : mais elevada no quarto superior ( 81,2 \% ) se comparada aos demais quartos de desempenho.
Observam -se diferenças estatisticamente significativas entre todos os quartos de desempenho. 4.4 SUFICIÊNCIA DAS INFORMAÇÕES/INSTRUÇÕES FORNECIDAS Ao avaliarem as informações/instruções fornecidas para a resolução das questões ( Questão 6 ), 82,2 \% dos respondentes da Área de Tecnologia em Logística de todo o Brasil afirmaram que estas eram até excessivas ou “ suficientes ” em todas ou na maioria das questões ( Gráfico 4.11, Gráfico 4.12, e, no Anexo II, a Tabela II.6 ).
Quanto à distribuição de respondentes pelas Grandes Regiões observa -se que a proporção de estudantes que consideraram as informações/instruções fornecidas até excessivas ou “ suficientes ” em todas ou na maioria das questões foi sempre superior a 78 \%, chegando a 88,6 \% na região Centro-Oeste.
As proporções das regiões Sul ( 85,0 \% ) e Centro-Oeste ( 88,6 \% ).
As proporções das regiões Sul ( 85,0 \% ) e Centro-Oeste ( 88,6 \% ) apresentam diferenças estatisticamente significativas em relação à região Norte ( 78,2 \% ) e Sudeste ( 81,2 \% ).
O percentual de participantes que avaliaram “ as informações/instruções ” como até excessivas ou “ suficientes ” em todas ou na maioria das questões foi mais elevado no quarto superior ( 88,9 \% ), percentual superior à média nacional ( 82,2 \% ).
Já no quarto inferior, a suficiência das informações/instruções declarada como até excessiva, em todas ou na maioria das questões foi percebida por 76,1 \% dos respondentes.
Essas proporções são crescentes com os quartos de desempenho.
Levando -se em conta o desempenho dos participantes, só não se observa diferença estatisticamente significativa entre o primeiro e o segundo quartos de desempenho, como mostra o Gráfico 4.12. 4.5 DIFICULDADE ENCONTRADA AO RESPONDER À PROVA Perguntados sobre as dificuldades com as quais se depararam ao responder à prova ( Questão 7 ), 10,5 \% dos estudantes apontaram o Desconhecimento do conteúdo.
Para 52,8 \%, a Forma diferente de abordagem do conteúdo foi indicada como dificuldade.
Já a Falta de motivação para fazer a prova foi a dificuldade apontada por 9,4 \% dos respondentes.
Considerando -se todo o Brasil, 22,7 \% dos respondentes afirmaram que não tiveram qualquer tipo de dificuldade para responder à prova ( Tabela II.7 no Anexo II ).
Os Gráficos 4.13 e 4.14 apresentam os percentuais de estudantes que apontaram o Desconhecimento do conteúdo como dificuldade percebida ao responder à prova.
Na análise por Grandes Regiões, o percentual de inscritos e presentes que apontaram o Desconhecimento do conteúdo como dificuldade ao responder à prova não superou 13 \%.
Os percentuais variaram de 6,9 \% na região Nordeste a 12,8 \% na Norte, sendo essa diferença significativa estatisticamente.
A Forma diferente de abordagem do conteúdo foi a escolha modal dos estudantes, com percentuais que variaram de 46,4 \% ( região Nordeste ) a 54,6 \% ( Sudeste ).
O percentual de alunos que citaram a Falta de motivação como dificuldade variou de 7,5 \% ( região Centro-Oeste ) a 10,2 \% ( região Sul ).
Os que declararam não ter qualquer tipo de dificuldade para responder à prova variaram de 19,7 \% na região Norte a 31,0 \% na Nordeste.
Com relação aos quartos de desempenho, o Desconhecimento do conteúdo foi a opção escolhida por 10,3 \% dos estudantes do quarto superior e 11,4 \% do quarto inferior.
Não há diferenças estatisticamente significativas entre os quartos.
A alternativa modal para os alunos, quando agregados pelos quartos de desempenho, foi que a dificuldade encontrada foi causada pela Forma diferente de abordagem do conteúdo : 50,6 \% no quarto inferior e 51,9 \% do quarto superior assim o responderam.
O percentual de alunos que citaram a Falta de motivação como dificuldade variou de 7,3 \% ( quarto superior ) a 13,4 \% ( quarto inferior ).
Os que declararam não ter qualquer tipo de dificuldade para responder à prova variaram de 19,4 \% no quarto inferior a 25,8 \% no superior. 4.6 CONTEÚDOS DAS QUESTÕES OBJETIVAS DA PROVA Ao analisarem os conteúdos das questões objetivas da prova ( Questão 8 ), um percentual muito pequeno dos estudantes avaliados, apenas 5,3 \%, afirmou que não estudou ainda a maioria desses conteúdos ( Gráficos 4.15, Gráfico 4.16, e a Tabelas II.8 no Anexo II ).
A maioria ( 71,1 \% ) afirmou ter estudado e aprendido muitos ou todos os conteúdos.
Na análise por Grande Região, a proporção de respondentes que escolheram a opção não estudou ainda a maioria desses conteúdos foi pequena, variando de 4,2 \% ( Nordeste ) a 9,9 \% ( Norte ).
Observam -se diferenças estatisticamente significativas entre a região Norte ( 9,9 \% ) e as regiões Nordeste, Sudeste e Sul.
Em todas as regiões, a maioria dos presentes afirmou ter estudado e aprendido muitos ou todos os conteúdos, com proporções variando entre 59,8 \% na região Norte e 76,4 \% na Centro-Oeste.
Considerando -se separadamente as opiniões de estudantes dos quatro quartos de desempenho, observa -se que, no quarto inferior, 9,1 \% ofereceram como resposta que não estudou ainda a maioria desses conteúdos, sendo 2,5 \% os do quarto superior com a mesma resposta.
Há diferenças estatisticamente significativas entre o quarto inferior e os demais, assim como entre o segundo quarto e o quarto superior.
Tendo em conta o quarto superior, 83,3 \% dos alunos afirmaram ter estudado e aprendido muitos ou todos os conteúdos.
No outro extremo, no primeiro quarto, 58,8 \% optaram pelas mesmas categorias. 4.7 TEMPO GASTO PARA CONCLUIR A PROVA Ao responderem sobre o tempo de conclusão da prova ( Questão 9 ), mais de três quartos dos estudantes ( 75,5 \% ) afirmou ter gasto Entre duas e quatro horas ( Gráfico 4.17, Gráfico 4.18 e, no Anexo II, a Tabela II.9 ).
Considerando -se as cinco Grandes Regiões brasileiras, o percentual dos que utilizaram Entre duas e quatro horas para finalizar a prova apenas na região Sul ( 72,5 \% ) foi inferior ao percentual nacional.
Nas demais Grandes Regiões, o percentual de alunos que dispensaram Entre duas e quatro horas para concluir a prova variou de 75,5 \% ( Sudeste ) a 82,5 \% ( Centro-Oeste ).
Observam -se diferenças estatisticamente significativas das regiões Sudeste e Sul em relação ao Nordeste ( 80,1 \% ) e também em relação ao Centro-Oeste, como mostra o Gráfico 4.17.
Uma vez sendo analisadas as alternativas escolhidas pelos estudantes que se situam nos diferentes quartos de desempenho, observa -se uma tendência crescente : uma proporção maior de participantes no quarto superior ( 83,2 \% ) declarou ter gasto Entre duas e quatro horas “ para concluir a prova ” quando comparadas com os dos quartos inferiores.
Há diferenças estatisticamente significativas entre todos os quartos de desempenho.

% matched lemmas: 
\end{textsample}
